\chapter{Dynamic Memory and ADTs in C}

\textbf{Readings:} CP:AMA 17.1--17.4, 19.

\section{Motivation}

So far, the size of every piece of data had to be known when the program was written: arrays
have a fixed length, and local variables disappear when their function returns. That causes
problems:
\begin{itemize}
  \item We cannot read an arbitrary number of values (how big should the array be?).
  \item A function cannot create data that outlives its own stack frame (returning a pointer to a
    local is a dangling pointer).
  \item An opaque structure cannot be created by a client, since the client does not know its
    size.
\end{itemize}
The solution is \emph{dynamic memory}: memory requested while the program is running, from the
\textbf{heap} section.

\section{Duration (Lifetime)}

\begin{center}
\begin{tabular}{lll}
  \toprule
  Kind & Where & Lifetime \\
  \midrule
  global/static & global or read-only data & the entire program \\
  local (automatic) & stack & until its block/function ends \\
  dynamic & heap & from \code{malloc} until \code{free} \\
  \bottomrule
\end{tabular}
\end{center}

\section{\code{malloc}}

\code{malloc} (``memory allocate'', in \code{<stdlib.h>}) requests a block of memory from the
heap:
\begin{lstlisting}[style=c]
void *malloc(size_t s);
\end{lstlisting}
It returns a pointer to the start of a block of (at least) \code{s} bytes, or \code{NULL} if
there is not enough memory. The return type \code{void *} is a \emph{generic pointer}: it can be
assigned to any pointer type without a cast.

\begin{lstlisting}[style=c]
#include <stdlib.h>

int *p = malloc(sizeof(int));          // one int
*p = 42;

int *arr = malloc(100 * sizeof(int));  // an array of 100 ints
arr[99] = 7;

struct posn *pp = malloc(sizeof(struct posn));
pp->x = 1;
\end{lstlisting}

\begin{important}[Using \code{malloc} correctly]
\begin{itemize}
  \item Always use \code{sizeof} --- never hard-code sizes like 4.
  \item The memory is \textbf{uninitialized} (garbage). Initialize it yourself, or use
    \code{calloc(n, size)}, which zero-fills.
  \item \code{malloc} can return \code{NULL}. In CS 136 we may assume it succeeds, but robust code
    checks.
  \item Every block that is \code{malloc}ed must eventually be \code{free}d exactly once.
\end{itemize}
\end{important}

\begin{center}
\begin{tikzpicture}[cell/.style={draw,minimum width=1.5cm,minimum height=0.7cm},
  ptr/.style={-{Stealth[length=2mm]},thick}]
  \node[font=\bfseries] at (0,1.2) {stack};
  \node[cell] (p) at (0,0) {0x9000};
  \node[left=1mm of p] {\code{p}};
  \node[font=\bfseries] at (5,1.2) {heap};
  \node[cell] (h) at (5,0) {42};
  \node[below=0mm of h,font=\scriptsize] {0x9000};
  \draw[ptr] (p.east) -- (h.west);
\end{tikzpicture}
\end{center}

\section{\code{free}}

\begin{lstlisting}[style=c]
void free(void *p);
\end{lstlisting}
\code{free(p)} returns the block that \code{p} points at to the heap so it can be reused.
\code{p} must be a pointer that was returned by \code{malloc} (or \code{calloc}/\code{realloc})
and not yet freed. \code{free(NULL)} does nothing.

\subsection{Memory errors}

\begin{itemize}
  \item \textbf{Memory leak:} losing the last pointer to a block without freeing it. The memory
    cannot be used again until the program ends. Long-running programs that leak eventually run
    out of memory. Seashell reports leaks when the program exits; leaks lose marks.
\begin{lstlisting}[style=c]
int *p = malloc(sizeof(int));
p = malloc(sizeof(int));     // leak: the first block is unreachable
\end{lstlisting}
  \item \textbf{Dangling pointer (use after free):} using a pointer after its block has been
    freed. The memory may have been reused for something else.
\begin{lstlisting}[style=c]
free(p);
*p = 5;          // undefined behaviour
p = NULL;        // good habit: then *p crashes immediately instead
\end{lstlisting}
  \item \textbf{Double free:} freeing the same block twice (often through two aliases): undefined
    behaviour, typically corrupting the heap.
  \item Freeing memory that was not \code{malloc}ed (e.g.\ \code{free(\&x)} on a local) is also
    undefined behaviour.
\end{itemize}

\begin{aside}
The heap does not have to be used in order: blocks can be freed in any order, which leaves
``holes'' (\emph{fragmentation}). The heap manager keeps track of free blocks. On Linux,
\code{valgrind} is a common tool to detect leaks and invalid accesses.
\end{aside}

\section{Documenting Dynamic Memory}

If a function returns memory that the caller must free, that is part of the interface and must be
documented:
\begin{lstlisting}[style=c]
// make_range(n) produces a new array {0, 1, ..., n-1}
// requires: n > 0
// effects: allocates memory (caller must free)
// time: O(n)
int *make_range(int n) {
  assert(n > 0);
  int *a = malloc(n * sizeof(int));
  for (int i = 0; i < n; ++i) {
    a[i] = i;
  }
  return a;
}

int main(void) {
  int *r = make_range(10);
  // ... use r ...
  free(r);
}
\end{lstlisting}
Unlike returning the address of a local, returning a heap pointer is fine: the block survives
until it is freed.

\section{\code{realloc}}

\begin{lstlisting}[style=c]
void *realloc(void *p, size_t newsize);
\end{lstlisting}
\code{realloc} resizes a previously allocated block to \code{newsize} bytes and returns a pointer
to the (possibly moved) block. The old contents are preserved (up to the smaller of the two
sizes). If it moves the block, it copies the data and frees the old block, so the old pointer
must no longer be used.
\begin{itemize}
  \item \code{realloc(NULL, s)} behaves like \code{malloc(s)}.
  \item Always use the returned pointer: \code{a = realloc(a, new\_len * sizeof(int));}
  \item In the worst case realloc must copy the whole block, so it is $\bigO(n)$ where $n$ is the
    size of the block.
\end{itemize}

\subsection{Reading a string of unknown length}

\begin{lstlisting}[style=c]
// read_line() reads a line of input (without the '\n') into a new string
//   or produces NULL at EOF
// effects: reads input, allocates memory (caller must free)
// time: O(n), where n is the length of the line
char *read_line(void) {
  int maxlen = 1;
  int len = 0;
  char *str = malloc(maxlen * sizeof(char));
  int c = getchar();
  if (c == EOF) {
    free(str);
    return NULL;
  }
  while (c != EOF && c != '\n') {
    if (len + 1 == maxlen) {        // need room for c and '\0'
      maxlen *= 2;
      str = realloc(str, maxlen * sizeof(char));
    }
    str[len] = c;
    ++len;
    c = getchar();
  }
  str[len] = '\0';
  return str;
}
\end{lstlisting}

\section{Dynamic Arrays and Amortized Analysis}

A \emph{dynamic array} (resizable array) keeps a pointer to a heap array, together with the
current length and the allocated capacity:
\begin{lstlisting}[style=c]
struct dynarray {
  int *data;
  int len;       // number of elements in use
  int maxlen;    // allocated capacity
};
\end{lstlisting}
When it is full and we want to append, we must grow it. How much should it grow?

\paragraph{Grow by one.} Each append calls \code{realloc}, which may copy all $k$ existing
elements. Appending $n$ elements costs $1 + 2 + \dots + n \in \bigO(n^2)$ --- $\bigO(n)$ per append.

\paragraph{Double the capacity.} When full, double \code{maxlen}. Most appends are $\bigO(1)$; an
expensive copy happens only when the length reaches $1, 2, 4, 8, \dots$. Appending $n$ elements
copies at most
\[ 1 + 2 + 4 + \dots + 2^{\floor{\log_2 n}} < 2n \]
elements in total, plus $\bigO(1)$ work per append, so the total is $\bigO(n)$.

\begin{definition}
The \emph{amortized} running time of an operation is the total time of a sequence of $n$
operations divided by $n$. With doubling, append is $\bigO(1)$ \emph{amortized}, even though a
single append can take $\bigO(n)$ time.
\end{definition}

\begin{important}[Doubling strategy]
Grow a dynamic array by a constant \emph{factor} (doubling), not by a constant \emph{amount}.
Doubling makes $n$ appends $\bigO(n)$ in total ($\bigO(1)$ amortized each), at the cost of up to
half the allocated memory being unused.
\end{important}

\begin{aside}
Shrinking works similarly: halve the capacity when the array becomes only a quarter full (not
half --- otherwise alternating append/remove at the boundary would resize every time).
\end{aside}

\section{ADTs in C}

Now we can build proper ADTs in C, combining:
\begin{itemize}
  \item an \textbf{opaque structure} declared in the header (clients can only hold pointers);
  \item \textbf{create} and \textbf{destroy} functions that allocate and free the structure;
  \item interface functions that take a pointer to the structure.
\end{itemize}

\subsection{Interface: \code{stack.h}}
\begin{lstlisting}[style=c]
// stack.h: a stack ADT of integers
#include <stdbool.h>

struct stack;   // opaque

// stack_create() creates a new empty stack
// effects: allocates memory (client must call stack_destroy)
// time: O(1)
struct stack *stack_create(void);

// stack_destroy(s) frees all memory used by s
// effects: s is no longer valid
// time: O(1)
void stack_destroy(struct stack *s);

// time: O(1)
bool stack_is_empty(const struct stack *s);

// requires: s is not empty
// time: O(1)
int stack_top(const struct stack *s);

// stack_pop(s) removes and produces the top item
// requires: s is not empty
// effects: modifies s
// time: O(1)
int stack_pop(struct stack *s);

// effects: modifies s
// time: O(1) amortized
void stack_push(int item, struct stack *s);
\end{lstlisting}

\subsection{Implementation: \code{stack.c}}
\begin{lstlisting}[style=c]
#include <assert.h>
#include <stdlib.h>
#include "stack.h"

struct stack {
  int len;
  int maxlen;
  int *data;
};

struct stack *stack_create(void) {
  struct stack *s = malloc(sizeof(struct stack));
  s->len = 0;
  s->maxlen = 1;
  s->data = malloc(s->maxlen * sizeof(int));
  return s;
}

void stack_destroy(struct stack *s) {
  assert(s);
  free(s->data);     // free the inner array first...
  free(s);           // ...then the structure itself
}

bool stack_is_empty(const struct stack *s) {
  assert(s);
  return s->len == 0;
}

int stack_top(const struct stack *s) {
  assert(s);
  assert(s->len > 0);
  return s->data[s->len - 1];
}

int stack_pop(struct stack *s) {
  assert(s);
  assert(s->len > 0);
  --s->len;
  return s->data[s->len];
}

void stack_push(int item, struct stack *s) {
  assert(s);
  if (s->len == s->maxlen) {
    s->maxlen *= 2;
    s->data = realloc(s->data, s->maxlen * sizeof(int));
  }
  s->data[s->len] = item;
  ++s->len;
}
\end{lstlisting}

\subsection{Client}
\begin{lstlisting}[style=c]
#include <stdio.h>
#include "stack.h"

int main(void) {
  struct stack *s = stack_create();
  for (int i = 0; i < 1000; ++i) {
    stack_push(i, s);              // no fixed maximum any more
  }
  while (!stack_is_empty(s)) {
    printf("%d\n", stack_pop(s));
  }
  stack_destroy(s);
}
\end{lstlisting}
The client cannot write \code{s->len} or \code{struct stack x;} because the structure is not
defined in \code{stack.h}; the compiler enforces the abstraction. We can later replace the
dynamic array by a linked list (Unit~11) without changing any client.

\begin{important}[Freeing nested structures]
Free memory ``inside out'': free what a structure points to before freeing the structure itself.
Freeing \code{s} first would make \code{s->data} a dangling access and leak the array.
\end{important}

\section{Summary}
\begin{itemize}
  \item The heap holds dynamic memory whose lifetime is controlled by the programmer.
  \item \code{malloc(n * sizeof(T))} allocates (uninitialized); \code{free} releases exactly once;
    \code{realloc} resizes (and may move, $\bigO(n)$).
  \item Errors: leaks, dangling pointers, double frees.
  \item Document ``allocates memory (caller must free)'' in the interface.
  \item Dynamic arrays with doubling give $\bigO(1)$ amortized append.
  \item C ADTs: opaque struct + create/destroy + interface functions taking pointers.
\end{itemize}
